\clearpage\SourceLeaf{316}

\begin{center}
[II.]

[DIE PHYSISCHE MÖGLICHKEIT DER DATEN IN POTHENOTS AUFGABE.]

[1.]
\end{center}

\begin{center}
\textsc{Gauss} an \textsc{Schumacher}. Göttingen, 13. April 1836.
\end{center}

\indent \mbox{. . . . . . .} Bei dieser Gelegenheit will ich doch auch einer
andern Aufgabe erwähnen, nemlich der sogenannten \textsc{Pothenots}chen.
Wenn drei Punkte $A$, $B$, $C$ gegeben sind, und zugleich die
Winkel, die $XA$, $XB$, $XC$ unter einander machen
(natürlich inclusive des \emph{sens}, d.\,i. ob die Winkel von der linken
nach der rechten [Seite] oder umgekehrt gezählt sind), so ist
$X$ (allgemein zu reden, d.\,i. mit Ausnahme des Falls, wo $A$,
$B$, $C$, $X$ auf Einem Kreise liegen) \textls[120]{bestimmt}. Gründen sich die
Data auf einen wirklichen concreten Fall, so versteht sich von selbst,
dass die Aufgabe, weil wirklich, auch möglich ist. Aber wenn jemand (um
Sie aufs Eis zu führen) Ihnen bloss die Zahlen aufgibt, so kann sich
hinterher ergeben, dass die Aufgabe physisch unmöglich ist, und dass Sie
statt einer der drei Richtungen die entgegengesetzte erhalten
($180^0$ verschieden von der aufgegebenen).

Die Frage ist nun, wie man \textls[120]{am einfachsten} die Bedingung der physischen
Möglichkeit durch eine Gleichung \textls[120]{zwischen den Datis} ausdrücken kann? Ich
habe, mich verwundernd, dass bei einer Aufgabe, worüber so viel
geschrieben ist, doch von allen, so viel ich weiss, ein so wesentlicher
Umstand ganz übersehen ist, — die Auflösung oft in meinen Vorlesungen
vorgetragen. Aber so ganz leicht muss sie doch wohl nicht zu finden sein;
wenigstens habe ich einmal einem sonst sehr geschickten jungen Mann (dem
sie von mir \textls[120]{nicht} vorgetragen war) sie als Thema einer Doctordissertation
vorgeschlagen, und der hatte, obgleich er sich lange damit gequält, die
Auflösung nicht finden können. \mbox{. . . . . . .}
